%! Unit 1, Lessons 1.0–1.7 — SAMPLE Test (Answer Key)
% The practice form of the Unit 1 test. Item for item it is the parallel form of ../actual_test/:
% the same 20 multiple-choice targets in the same order and the same five free-response
% questions with the same parts and points (MC 20 x 3 = 60, FRQ 5 x 8 = 40, total 100;
% 25 questions) --- every context different, so working this one teaches the moves, not the answers.
% Merged into the unit student packet via ../../sample_test/ (see shared/tests.mk).
% NAME ROW: \namedateperiod IS carried --- a unit test is taken in a testing setting.
% NO SKETCH ITEMS: every display is pre-drawn in TikZ and read.
%
% BLUEPRINT
%   MC 1 individuals                      bike-repair work orders          1.0
%   MC 2 statistical question             bike-repair work orders          1.0
%   MC 3 categorical vs quantitative      bike-repair work orders (ZIP)    1.1
%   MC 4 grouping (binning)               gym workout minutes              1.1
%   MC 5 relative frequency               how 160 students watch movies    1.2
%   MC 6 groups of unequal size           two summer camps, swimming       1.2
%   MC 7 discrete vs continuous           mile time to the second          1.3
%   MC 8 read a stem plot (median)        15 April high temperatures       1.3
%   MC 9 interval width                   50 battery lifetimes             1.3
%   MC10 shape from a histogram           40 used-textbook prices          1.4
%   MC11 unusual features (dotplot)       20 song lengths                  1.4
%   MC12 mean vs median from shape        80 home sale prices              1.5
%   MC13 resistance (a typo)              9 lunch-special prices           1.5
%   MC14 standard deviation in context    30 dog-food bags                 1.5
%   MC15 percentile                       state reading test               1.5
%   MC16 1.5 x IQR rule                   40 concert ticket prices         1.6
%   MC17 every section holds a quarter    30 basketball games              1.6
%   MC18 parallel boxplots                two urgent-care clinics          1.6
%   MC19 parameter vs statistic           town residents' ages             1.7
%   MC20 z-scores across distributions    history vs chemistry tests       1.7
%   FRQ1 individuals, variables, SQ       city playgrounds                 1.0, 1.1, 1.3
%   FRQ2 relative frequency comparison    two high schools' senior plans   1.2
%   FRQ3 histogram, describe              50 trout lengths                 1.3, 1.4
%   FRQ4 five numbers, outlier, center    15 used video-game prices        1.5, 1.6
%   FRQ5 empirical rule, z-scores         egg weights                      1.7
% ANSWERS (MC): 1B 2D 3B 4E 5C 6A 7D 8B 9E 10C 11A 12D 13C 14E 15A 16D 17B 18E 19C 20A
%
% THE DATA (every number checked)
%  MC5  92/36/20/12 = 160; theater 36/160 = 0.225.
%  MC6  Camp A 54/120 = 0.45; Camp B 39/75 = 0.52.
%  MC8  54 56 58 59 60 61 63 63 65 67 68 70 72 72 76; median (8th) 63; mean 964/15 = 64.3.
%  MC10 bins of $10 from 0: 12 11 7 5 3 1 1 (40), skewed right.
%  MC11 2.5:1 3.0:3 3.5:6 4.0:5 4.5:3 5.0:1 8.0:1 (20).
%  MC13 6 7 7 8 9 9 10 11 12 -> 120: Q1 7, med 9, Q3 10.5 both ways; mean, range move.
%  MC16 18/30/38/46/95: IQR 16, fences 6 and 70; 95 > 70.
%  MC17 52/60/64/75/88.
%  MC18 A 10/18/24/30/45 (IQR 12); B 15/26/30/34/40 (IQR 8).
%  MC20 Maya (84-76)/5 = 1.6; Jordan (88-80)/8 = 1.0.
%  FRQ2 North 120/55/50/25 (250) = 48/22/20/10%; South 88/32/28/12 (160) = 55/20/17.5/7.5%.
%  FRQ3 bins of 4 cm from 12: 2 0 3 5 9 15 11 5 (50); cum 2 2 5 10 19 34 45 50.
%  FRQ4 8 10 12 12 14 15 15 16 18 20 21 22 24 25 58: Q1 12, med 16, Q3 22, IQR 10,
%       fences -3 and 37; sum 290, mean 19.33; 9 prices below the mean.
%  FRQ5 N(58, 4): 50-66 95%; >62 16%; 51 -> z -1.75; 65 -> 1.75 vs farm 2 N(64, 6) 73 -> 1.5.
%
% LOCKSTEP: the blank ../../tests/practice_test/ is generated FROM THIS FILE (edit here, regenerate).
% Every ans mark on an MC option prints bare there; every other ans mark becomes a phantom of
% itself; every run of n anslines becomes writelines{n}; work blocks stay byte-identical.
\documentclass[12pt]{article}
\usepackage{apstats-article}
\usepackage{apstats-key}   % pulls in -boxes; do NOT also load -boxes

\begin{document}

\pageheader{Unit 1 \quad Lessons 1.0--1.7}{Sample Test --- Answer Key}

\namedateperiod

{\small\hfill Score:\ \blank{1.4cm}\ / 100}

\vspace{0.06in}

\boxguard[8]
\begin{remindbox}
\small
This is the \textbf{practice} test --- not scored. Work it in one period, notes closed, then check
it against the key. It matches the real test item for item; every answer must be in context.
\end{remindbox}

\vspace{0.10in}

\boxguard[16]
\begin{headlinebox}{goldbg}
\small
\textbf{Section I --- Multiple Choice} \hfill \textbf{60 points, 3 each}\\
Circle the single best answer. Read all five options --- more than one of them is about the
context, not just the arithmetic.
\end{headlinebox}

\vspace{0.10in}

\boxguard[10]
\noindent\textbf{Questions 1--3} use these records from a bicycle repair shop. Each row is one
bike brought in for repair.

\vspace{4pt}
{\centering\small
\renewcommand{\arraystretch}{1.15}
\begin{tabular}{@{} c l c c c @{}}
\textbf{Order \#} & \textbf{Bike type} & \textbf{Repair time (min)} & \textbf{Cost (\$)} &
\textbf{Customer ZIP} \\ \hline
4417 & Road     & 45 & 62.50  & 30214 \\
4418 & Mountain & 90 & 118.00 & 30217 \\
4419 & Kids'    & 20 & 25.00  & 30214 \\
4420 & Hybrid   & 35 & 48.75  & 30220 \\
\end{tabular}\par}

\vspace{0.06in}

\begin{enumerate}[leftmargin=1.9em, itemsep=0.13in]

  \boxguard[8]
  \item In the shop's records, the \textbf{individuals} are
        \begin{enumerate}[label=\textbf{(\Alph*)}, leftmargin=2.2em, itemsep=2pt]
          \item the bike types.
          \item \ans{the bikes brought in for repair.}
          \item the customers' ZIP codes.
          \item the repair times.
          \item the bicycle repair shop.
        \end{enumerate}

  \boxguard[8]
  \item Which of the following is a \textbf{statistical question} about the shop?
        \begin{enumerate}[label=\textbf{(\Alph*)}, leftmargin=2.2em, itemsep=2pt]
          \item How long did the repair on order 4417 take?
          \item What type of bike is order 4420?
          \item Is the shop open on Sundays?
          \item \ans{How much do repair times vary across the shop's orders last month?}
          \item How many mechanics work at the shop?
        \end{enumerate}

  \boxguard[8]
  \item How many of the five columns are \textbf{quantitative} variables?
        \begin{enumerate}[label=\textbf{(\Alph*)}, leftmargin=2.2em, itemsep=2pt]
          \item One
          \item \ans{Two}
          \item Three
          \item Four
          \item Five
        \end{enumerate}

  \boxguard[11]
  \item A gym records each member's weekly workout time in minutes, then replaces it with a
        label: \emph{Light} (under 90 min), \emph{Moderate} (90 to 180 min), or \emph{Heavy}
        (over 180 min). Which statement is true?
        \begin{enumerate}[label=\textbf{(\Alph*)}, leftmargin=2.2em, itemsep=2pt]
          \item The new variable is still quantitative, because the labels are based on minutes.
          \item The new variable is quantitative, because the labels are in order.
          \item Nothing is lost: the minutes can be recovered from the labels.
          \item The median workout time can still be found exactly from the labels.
          \item \ans{The new variable is categorical, and the mean workout time can no longer be
                found from it.}
        \end{enumerate}

  \boxguard[12]
  \item \begin{minipage}[t]{0.50\linewidth}\raggedright
        A survey asked 160 high-school students how they most often watch movies. What is the
        \textbf{relative frequency} of \emph{Theater}?
        \end{minipage}\hfill
        \begin{minipage}[t]{0.44\linewidth}
        \centering\small
        \renewcommand{\arraystretch}{1.1}
        \begin{tabular}[t]{@{} l c @{}}
        \textbf{Way to watch} & \textbf{Frequency} \\ \hline
        Streaming & 92 \\
        Theater   & 36 \\
        Cable TV  & 20 \\
        DVD       & 12 \\ \hline
        \textbf{Total} & 160 \\
        \end{tabular}
        \end{minipage}
        \begin{enumerate}[label=\textbf{(\Alph*)}, leftmargin=2.2em, itemsep=2pt]
          \item $0.036$
          \item $0.180$
          \item \ans{$0.225$}
          \item $0.360$
          \item $0.575$
        \end{enumerate}

  \boxguard[11]
  \item At Camp Alder, 54 of its 120 campers named swimming as their favorite activity. At
        Camp Birch, 39 of its 75 campers did. Which statement is best supported?
        \begin{enumerate}[label=\textbf{(\Alph*)}, leftmargin=2.2em, itemsep=2pt]
          \item \ans{Swimming is more popular at Camp Birch, because 52\% is greater than 45\%.}
          \item Swimming is more popular at Camp Alder, because 54 is greater than 39.
          \item Swimming is equally popular at the two camps.
          \item No comparison is possible, because the camps have different numbers of campers.
          \item Swimming is more popular at Camp Alder, because it has more campers.
        \end{enumerate}

  \boxguard[9]
  \item Which of these variables is \textbf{continuous}?
        \begin{enumerate}[label=\textbf{(\Alph*)}, leftmargin=2.2em, itemsep=2pt]
          \item The number of siblings a student has
          \item The number of text messages a student sent today
          \item The number of pets in a household
          \item \ans{The time a student takes to run a mile, recorded to the nearest second}
          \item The number of songs on a playlist
        \end{enumerate}

  \boxguard[12]
  \item \begin{minipage}[t]{0.52\linewidth}\raggedright
        The stem plot shows the daily high temperature ($^\circ$F) in a city on 15 days in
        April. What is the \textbf{median} daily high temperature?
        \end{minipage}\hfill
        \begin{minipage}[t]{0.42\linewidth}
        \centering\small
        \begin{tabular}[t]{@{} r | l @{}}
        5 & 4 6 8 9 \\
        6 & 0 1 3 3 5 7 8 \\
        7 & 0 2 2 6 \\
        \end{tabular}\par\vspace{2pt}
        {\footnotesize Key: $5\,|\,8 = 58^\circ$F}
        \end{minipage}
        \begin{enumerate}[label=\textbf{(\Alph*)}, leftmargin=2.2em, itemsep=2pt]
          \item $61^\circ$F
          \item \ans{$63^\circ$F}
          \item $64.3^\circ$F
          \item $65^\circ$F
          \item $7^\circ$F
        \end{enumerate}

  \boxguard[12]
  \item A student made two histograms of the \emph{same} 50 battery lifetimes (hours). With
        intervals 1 hour wide, the histogram shows a gap between 12 and 14 hours. With intervals
        4 hours wide, no gap appears. Which statement is correct?
        \begin{enumerate}[label=\textbf{(\Alph*)}, leftmargin=2.2em, itemsep=2pt]
          \item One of the two histograms must have been drawn incorrectly.
          \item The wider intervals changed some of the battery lifetimes.
          \item Only the histogram with more bars is a true histogram.
          \item The gap shows that the data set contains an outlier.
          \item \ans{Both are correct displays of the same data; changing the interval width
                changes what the display shows without changing any value.}
        \end{enumerate}

  \boxguard[14]
  \item \begin{minipage}[t]{0.46\linewidth}\raggedright
        The histogram shows the prices of 40 used textbooks sold online. Which best describes
        the \textbf{shape} of the distribution?
        \end{minipage}\hfill
        \begin{minipage}[t]{0.50\linewidth}
        \centering
        \begin{tikzpicture}[x=0.075cm, y=0.22cm, baseline=(current bounding box.north)]
          \foreach \y in {0,4,8,12} {
            \draw[linegray, very thin] (0,\y) -- (72,\y);
            \node[left, font=\tiny] at (0,\y) {\y};}
          \foreach \l/\c in {0/12, 10/11, 20/7, 30/5, 40/3, 50/1, 60/1}
            \draw[fill=skymid] (\l,0) rectangle (\l+10,\c);
          \draw[thick] (0,0) -- (73,0);
          \draw[thick] (0,0) -- (0,13);
          \foreach \x in {0,10,...,70} { \node[below, font=\tiny] at (\x,0) {\x}; }
          \node[font=\tiny] at (36,-3) {Price (dollars)};
          \node[rotate=90, font=\tiny] at (-9,6.5) {Count};
        \end{tikzpicture}
        \end{minipage}
        \begin{enumerate}[label=\textbf{(\Alph*)}, leftmargin=2.2em, itemsep=2pt]
          \item Skewed left, because most of the prices pile up on the left.
          \item Roughly symmetric, because there is only one peak.
          \item \ans{Skewed right, because a few prices are much higher than the rest.}
          \item Uniform, because every interval has at least one textbook.
          \item Bimodal, because the two lowest intervals are both tall.
        \end{enumerate}

  \boxguard[14]
  \item \begin{minipage}[t]{0.40\linewidth}\raggedright
        The dotplot shows the lengths, in minutes, of the 20 songs on a playlist. Which
        statement is best?
        \end{minipage}\hfill
        \begin{minipage}[t]{0.56\linewidth}
        \centering
        \begin{tikzpicture}[x=1.15cm, y=0.22cm, baseline=(current bounding box.north)]
          \draw[thick] (2,0) -- (8.5,0);
          \foreach \x in {2,3,...,8} {
            \draw (\x,-0.4) -- (\x,0.4); \node[below, font=\tiny] at (\x,-0.3) {\x};}
          \foreach \x/\n in {2.5/1, 3/3, 3.5/6, 4/5, 4.5/3, 5/1, 8/1}
            \foreach \k in {1,...,\n} \fill[navy] (\x,\k*0.9) circle (2.2pt);
          \node[font=\tiny] at (5.25,-3) {Song length (minutes)};
        \end{tikzpicture}
        \end{minipage}
        \begin{enumerate}[label=\textbf{(\Alph*)}, leftmargin=2.2em, itemsep=2pt]
          \item \ans{The 8-minute song is a possible outlier, set off by a gap from 5 to 8 minutes;
                it should be investigated, not deleted.}
          \item The 8-minute song is an outlier, so it should be deleted from the data.
          \item There are no unusual features, because most songs last 3 to 4.5 minutes.
          \item The distribution is bimodal, with peaks at 3.5 and 8 minutes.
          \item The gap from 2.5 to 3 minutes is the most unusual feature.
        \end{enumerate}

  \boxguard[11]
  \item The distribution of the 80 home sale prices in a town last year is strongly
        \textbf{skewed to the right}. Which statement is most likely true?
        \begin{enumerate}[label=\textbf{(\Alph*)}, leftmargin=2.2em, itemsep=2pt]
          \item The mean is less than the median, so the mean better describes a typical price.
          \item The mean and median are equal, because both measure center.
          \item The mean is greater than the median, so the mean better describes a typical
                price.
          \item \ans{The mean is greater than the median, so the median better describes a
                typical price.}
          \item Nothing can be said about the mean and median without the data.
        \end{enumerate}

  \boxguard[11]
  \item The prices of the 9 lunch specials at a diner are \$6, 7, 7, 8, 9, 9, 10, 11, and 12.
        The \$12 price is mistakenly entered as \$120. Which summaries change?
        \begin{enumerate}[label=\textbf{(\Alph*)}, leftmargin=2.2em, itemsep=2pt]
          \item Both the mean and the median change.
          \item The median changes, but the mean does not.
          \item \ans{The mean and the range change; the median and the IQR do not.}
          \item None of the summaries change, because only one price is wrong.
          \item Only the IQR changes.
        \end{enumerate}

  \boxguard[11]
  \item A pet-food company fills bags labeled 20 lb. For 30 bags, the standard deviation of the
        weights is 0.3 lb. Which is the best interpretation?
        \begin{enumerate}[label=\textbf{(\Alph*)}, leftmargin=2.2em, itemsep=2pt]
          \item Every bag weighs within 0.3 lb of the mean weight.
          \item A typical bag weighs 0.3 lb.
          \item The heaviest bag weighs 0.3 lb more than the lightest bag.
          \item About 30\% of the bags weigh more than the mean weight.
          \item \ans{The bag weights typically vary by about 0.3 lb from the mean weight.}
        \end{enumerate}

  \boxguard[10]
  \item A student's score on a state reading test is at the \textbf{82nd percentile}. Which
        statement is correct?
        \begin{enumerate}[label=\textbf{(\Alph*)}, leftmargin=2.2em, itemsep=2pt]
          \item \ans{About 82\% of the students who took the test scored at or below this
                student.}
          \item The student answered 82\% of the questions correctly.
          \item The student scored 82 points.
          \item About 18\% of the students scored at or below this student.
          \item The student scored higher than exactly 82 other students.
        \end{enumerate}

  \boxguard[12]
  \item The five-number summary of the ticket prices (dollars) for 40 concerts at a venue is
        \[ \text{Min} = 18 \quad Q_1 = 30 \quad \text{Median} = 38 \quad Q_3 = 46 \quad
           \text{Max} = 95. \]
        Which statement is true?
        \begin{enumerate}[label=\textbf{(\Alph*)}, leftmargin=2.2em, itemsep=2pt]
          \item Both \$18 and \$95 are outliers, because they are the minimum and maximum.
          \item There are no outliers, because the IQR of 16 is small.
          \item \$95 is an outlier, because it is more than 1.5 times the median.
          \item \ans{\$95 is an outlier, because it is greater than $Q_3 + 1.5 \times
                \text{IQR} = 70$.}
          \item Outliers cannot be identified from a five-number summary.
        \end{enumerate}

  \boxguard[14]
  \item \begin{minipage}[t]{0.46\linewidth}\raggedright
        The boxplot shows the points a basketball team scored in each of its 30 games. Which
        statement is true?
        \end{minipage}\hfill
        \begin{minipage}[t]{0.50\linewidth}
        \centering
        \begin{tikzpicture}[x=0.17cm, y=1cm, baseline=(current bounding box.north)]
          \draw[thick] (50,0) -- (90,0);
          \foreach \x in {50,55,...,90} {
            \draw (\x,-0.07) -- (\x,0.07); \node[below, font=\tiny] at (\x,-0.05) {\x};}
          \draw (52,0.65) -- (60,0.65);
          \draw (75,0.65) -- (88,0.65);
          \draw[fill=skymid] (60,0.4) rectangle (75,0.9);
          \draw[thick] (64,0.4) -- (64,0.9);
          \draw (52,0.5) -- (52,0.8);
          \draw (88,0.5) -- (88,0.8);
          \node[font=\tiny] at (70,-0.5) {Points scored};
        \end{tikzpicture}
        \end{minipage}
        \begin{enumerate}[label=\textbf{(\Alph*)}, leftmargin=2.2em, itemsep=2pt]
          \item More games had scores from 64 to 75 than from 60 to 64, because that part of the
                box is wider.
          \item \ans{About a quarter of the games had scores from 60 to 64, and about a quarter
                had scores from 64 to 75.}
          \item About half of the games had scores above 75.
          \item The team's mean score was 64 points.
          \item The range of the scores is 28 points.
        \end{enumerate}

  \boxguard[15]
  \item \begin{minipage}[t]{0.40\linewidth}\raggedright
        The parallel boxplots show patients' waiting times (minutes) at two urgent-care clinics.
        Which comparison is best?
        \end{minipage}\hfill
        \begin{minipage}[t]{0.56\linewidth}
        \centering
        \begin{tikzpicture}[x=0.15cm, y=1cm, baseline=(current bounding box.north)]
          \draw[thick] (5,0) -- (50,0);
          \foreach \x in {5,10,...,50} {
            \draw (\x,-0.07) -- (\x,0.07); \node[below, font=\tiny] at (\x,-0.05) {\x};}
          % Clinic A (upper)
          \draw (10,1.45) -- (18,1.45); \draw (30,1.45) -- (45,1.45);
          \draw[fill=skymid] (18,1.25) rectangle (30,1.65); \draw[thick] (24,1.25) -- (24,1.65);
          \draw (10,1.33) -- (10,1.57); \draw (45,1.33) -- (45,1.57);
          \node[left, font=\tiny] at (5,1.45) {A};
          % Clinic B (lower)
          \draw (15,0.6) -- (26,0.6); \draw (34,0.6) -- (40,0.6);
          \draw[fill=goldbg] (26,0.4) rectangle (34,0.8); \draw[thick] (30,0.4) -- (30,0.8);
          \draw (15,0.48) -- (15,0.72); \draw (40,0.48) -- (40,0.72);
          \node[left, font=\tiny] at (5,0.6) {B};
          \node[font=\tiny] at (27.5,-0.5) {Waiting time (minutes)};
        \end{tikzpicture}
        \end{minipage}
        \begin{enumerate}[label=\textbf{(\Alph*)}, leftmargin=2.2em, itemsep=2pt]
          \item Clinic A has the higher median wait, because its box reaches farther right.
          \item The clinics are equally variable, because both boxes are about the same height.
          \item Clinic A's waits are longer, because its maximum is larger.
          \item Every patient at Clinic B waited longer than every patient at Clinic A.
          \item \ans{Clinic B has the higher median wait (30 vs.\ 24 minutes) and the smaller IQR
                (8 vs.\ 12 minutes).}
        \end{enumerate}

  \boxguard[11]
  \item A town's records show that the mean age of all 4,200 residents is 38.6 years. A reporter
        surveys 50 residents and finds that their mean age is 41.2 years. Which is correct?
        \begin{enumerate}[label=\textbf{(\Alph*)}, leftmargin=2.2em, itemsep=2pt]
          \item 38.6 is a statistic $\bar x$, and 41.2 is a parameter $\mu$.
          \item Both numbers are parameters, because both describe residents.
          \item \ans{38.6 is a parameter $\mu$, and 41.2 is a statistic $\bar x$.}
          \item Both numbers are statistics, because both are means.
          \item 41.2 is a parameter, because it is the larger number.
        \end{enumerate}

  \boxguard[13]
  \item Maya scored 84 on a history test where the class mean was 76 and the standard deviation
        was 5. Jordan scored 88 on a chemistry test where the class mean was 80 and the standard
        deviation was 8. Who did better \emph{relative to their own class}?
        \begin{enumerate}[label=\textbf{(\Alph*)}, leftmargin=2.2em, itemsep=2pt]
          \item \ans{Maya, because Maya's $z$-score of 1.6 is greater than Jordan's 1.0.}
          \item Jordan, because 88 is greater than 84.
          \item They did equally well, because both scored 8 points above their class mean.
          \item Jordan, because Jordan's $z$-score of 1.6 is greater than Maya's 1.0.
          \item They cannot be compared, because the tests are in different subjects.
        \end{enumerate}

\end{enumerate}

\vspace{0.14in}

% No \newpage here: the guard keeps the Section II strip off a page edge.
\boxguard[30]
\begin{headlinebox}{goldbg}
\small
\textbf{Section II --- Free Response} \hfill \textbf{40 points, 8 each}\\
Show your reasoning. Answer \emph{in context} --- name the variables and their units.
\end{headlinebox}

\vspace{0.10in}

\begin{enumerate}[leftmargin=1.9em, itemsep=0.16in, resume]

  \boxguard[24]
  \item \textbf{(8 points)} A city's parks department keeps these records on its playgrounds.

        \vspace{4pt}
        {\centering\small
        \renewcommand{\arraystretch}{1.12}
        \begin{tabular}{@{} c l c c l @{}}
        \textbf{Park \#} & \textbf{Neighborhood} & \textbf{Swings} & \textbf{Area (sq ft)} &
        \textbf{Surface} \\ \hline
        104 & Eastgate   & 6 & 2,450 & Wood chips \\
        117 & Mill Creek & 4 & 1,820 & Rubber \\
        122 & Eastgate   & 8 & 3,100 & Rubber \\
        135 & Riverside  & 2 & 960   & Sand \\
        141 & Hillcrest  & 6 & 2,275 & Wood chips \\
        \end{tabular}\par}

        \vspace{6pt}
        \textbf{(a)} What are the individuals in this data set? Is \emph{Eastgate} a variable or
        a value? Explain.\\[4pt]
        \ansline{The individuals are the city's playgrounds --- one row each.}\\
        \ansline{\emph{Eastgate} is a value: one entry under the variable \emph{Neighborhood}.}

        \boxguard[12]
        \vspace{6pt}
        \textbf{(b)} Classify each variable. For each quantitative variable, also say whether it
        is discrete or continuous.

        \vspace{4pt}
        {\centering\small
        \renewcommand{\arraystretch}{1.5}
        \begin{tabular}{@{} l | p{4.2cm} | p{4.2cm} @{}}
        \textbf{Variable} & \textbf{Categorical or quantitative?} &
        \textbf{If quantitative: discrete or continuous?} \\ \hline
        Park \#      & \ans{Categorical}           & \\ \hline
        Neighborhood & \ans{Categorical}           & \\ \hline
        Swings       & \ans{Quantitative}          & \ans{Discrete} \\ \hline
        Area         & \ans{Quantitative}          & \ans{Continuous} \\ \hline
        Surface      & \ans{Categorical}           & \\
        \end{tabular}\par}

        \boxguard[6]
        \vspace{6pt}
        \textbf{(c)} A council member asks, ``How many swings does the playground at Park 122
        have?'' Is this a statistical question? Explain. Then write a statistical question these
        data could answer.\\[4pt]
        \ansline{No: it has one answer (8 swings), read from one row --- no variability.}\\
        \ansline{Statistical: ``How much does playground area vary across the city's}\\
        \ansline{playgrounds?'' --- it expects different answers from different playgrounds.}

  \boxguard[24]
  \item \textbf{(8 points)} Two high schools asked their seniors about their plans after
        graduation.

        \vspace{4pt}
        \begin{minipage}[c]{0.50\linewidth}
        \centering\small
        \setlength{\tabcolsep}{4pt}\renewcommand{\arraystretch}{1.15}
        \begin{tabular}{@{} l | c c @{}}
         & \textbf{North High} & \textbf{South High} \\ \hline
        4-year college & 120 & 88 \\
        2-year college & 55  & 32 \\
        Work           & 50  & 28 \\
        Military       & 25  & 12 \\ \hline
        \textbf{Total} & 250 & 160 \\
        \end{tabular}
        \end{minipage}\hfill
        \begin{minipage}[c]{0.46\linewidth}
        \centering
        \begin{tikzpicture}[x=0.62cm, y=0.022cm]
          \foreach \y in {0,40,80,120} {
            \draw[linegray, very thin] (0,\y) -- (8.4,\y);
            \node[left, font=\tiny] at (0,\y) {\y};}
          \foreach \i/\n/\s in {0/120/88, 1/55/32, 2/50/28, 3/25/12} {
            \fill[navy]    (0.3+2*\i,0) rectangle (1.0+2*\i,\n);
            \fill[goldacc] (1.0+2*\i,0) rectangle (1.7+2*\i,\s);}
          \draw[thick] (0,0) -- (8.4,0);
          \draw[thick] (0,0) -- (0,130);
          \foreach \i/\lab in {0/4-yr, 1/2-yr, 2/Work, 3/Military}
            \node[below, font=\tiny] at (1.0+2*\i,0) {\lab};
          \node[rotate=90, font=\tiny] at (-1.1,65) {Number of seniors};
          \fill[navy] (4.6,128) rectangle (4.9,138); \node[right, font=\tiny] at (4.9,133) {North};
          \fill[goldacc] (6.5,128) rectangle (6.8,138); \node[right, font=\tiny] at (6.8,133) {South};
        \end{tikzpicture}
        \end{minipage}

        \vspace{6pt}
        \textbf{(a)} What proportion of North High seniors plan to attend a 4-year college? What
        proportion of South High seniors?\\[4pt]
        \ansline{North: $120/250 = 0.48$ (48\%).}\\
        \ansline{South: $88/160 = 0.55$ (55\%).}

        \boxguard[6]
        \vspace{6pt}
        \textbf{(b)} A reporter wrote: ``North High sends more of its seniors to 4-year colleges
        --- just look how much taller its bar is.'' Explain what is wrong with this reasoning.\\[4pt]
        \ansline{The bars show counts, and North has 250 seniors to South's 160, so its bars}\\
        \ansline{are taller just because the school is bigger. Compare relative frequencies:}\\
        \ansline{48\% of North seniors vs.\ 55\% of South seniors --- South's share is larger.}

        \boxguard[6]
        \vspace{6pt}
        \textbf{(c)} Do \emph{most} North High seniors plan to attend a 4-year college? Do most
        South High seniors? Justify with numbers.\\[4pt]
        \ansline{North: no. 48\% is less than half --- a 4-year college is North's most common}\\
        \ansline{plan, but not a majority. South: yes. 55\% is more than half of the}\\
        \ansline{160 seniors.}

  \boxguard[26]
  \item \textbf{(8 points)} A biologist measured the lengths (cm) of 50 trout caught in a lake.
        The histogram shows the results.

        \vspace{4pt}
        {\centering
        \begin{tikzpicture}[x=0.27cm, y=0.24cm]
          \foreach \y in {0,4,8,12,16} {
            \draw[linegray, very thin] (10,\y) -- (46,\y);
            \node[left, font=\tiny] at (10,\y) {\y};}
          \foreach \l/\c in {12/2, 16/0, 20/3, 24/5, 28/9, 32/15, 36/11, 40/5}
            \draw[fill=skymid] (\l,0) rectangle (\l+4,\c);
          \draw[thick] (10,0) -- (46,0);
          \draw[thick] (10,0) -- (10,17);
          \foreach \x in {12,16,...,44} { \node[below, font=\tiny] at (\x,0) {\x}; }
          \node[font=\tiny] at (28,-2.6) {Length (cm)};
          \node[rotate=90, font=\tiny] at (7.6,8) {Number of trout};
        \end{tikzpicture}\par}

        \vspace{6pt}
        \textbf{(a)} How many trout were shorter than 28 cm? What percent of the catch is that?\\[4pt]
        \ansline{$2 + 0 + 3 + 5 = 10$ trout were shorter than 28 cm.}\\
        \ansline{$10/50 = 0.20$, so 20\% of the catch.}

        \boxguard[8]
        \vspace{6pt}
        \textbf{(b)} Describe the distribution of trout lengths.\\[4pt]
        \ansline{Shape: skewed left --- a long tail toward short lengths; one peak at 32--36 cm.}\\
        \ansline{Unusual: a gap at 16--20 cm; the 2 trout at 12--16 cm are possible outliers.}\\
        \ansline{Center: the median length is in the 32--36 cm interval (about 33 cm).}\\
        \ansline{Variability: lengths run from about 12 cm to 44 cm, a range of about 32 cm.}

        \boxguard[6]
        \vspace{6pt}
        \textbf{(c)} Can the exact median length be found from the histogram? If not, which
        interval contains it? Explain.\\[4pt]
        \ansline{No --- a histogram groups the values. 19 trout are under 32 cm and 34 are}\\
        \ansline{under 36 cm, so the 25th and 26th lengths (the median) are in 32--36 cm.}

  \boxguard[22]
  \item \textbf{(8 points)} A resale shop's prices (dollars) for 15 used video games, in order:
        \[ 8 \quad 10 \quad 12 \quad 12 \quad 14 \quad 15 \quad 15 \quad 16 \quad 18 \quad 20
           \quad 21 \quad 22 \quad 24 \quad 25 \quad 58 \]

        \textbf{(a)} Find the five-number summary.
        \begin{work}
          \text{Min} &= 8, \quad Q_1 = 12 \ (\text{median of the lower 7}) \\
          \text{Median} &= 16 \ (\text{the 8th price}) \\
          Q_3 &= 22 \ (\text{median of the upper 7}), \quad \text{Max} = 58
        \end{work}

        \boxguard[10]
        \vspace{6pt}
        \textbf{(b)} Use the $1.5 \times \text{IQR}$ rule to decide whether \$58 is an outlier. If
        it is, where does the upper whisker of a modified boxplot end?
        \begin{work}
          \text{IQR} &= 22 - 12 = 10 \\
          Q_3 + 1.5(\text{IQR}) &= 22 + 15 = 37
        \end{work}
        \ansline{\$58 $>$ \$37, so \$58 is an outlier. The upper whisker ends at \$25.}

        \boxguard[6]
        \vspace{6pt}
        \textbf{(c)} The mean price is \$19.33. Which better describes a typical price here, the
        mean or the median? Explain.\\[4pt]
        \ansline{The median, \$16. It is resistant: the \$58 outlier barely moves it. The mean}\\
        \ansline{is nonresistant --- the \$58 game drags it up to \$19.33, above 9 of the}\\
        \ansline{15 prices, so it is not a typical price.}

  \boxguard[20]
  \item \textbf{(8 points)} The weights of eggs from a farm are approximately normal, with mean
        $\mu = 58$ grams and standard deviation $\sigma = 4$ grams.

        \vspace{4pt}
        {\centering
        \begin{tikzpicture}[x=1.0cm, y=3.2cm]
          \draw[thick, navy, domain=-3.6:3.6, samples=80, smooth]
            plot (\x, {exp(-\x*\x/2)});
          \draw[thick] (-3.8,0) -- (3.8,0);
          \foreach \k/\lab in {-3/46, -2/50, -1/54, 0/58, 1/62, 2/66, 3/70} {
            \draw (\k,-0.04) -- (\k,0.04); \node[below, font=\tiny] at (\k,-0.02) {\lab};}
          \node[font=\tiny] at (0,-0.28) {Egg weight (grams)};
        \end{tikzpicture}\par}

        \vspace{4pt}
        \textbf{(a)} About what percent of the eggs weigh between 50 and 66 grams?\\[4pt]
        \ansline{50 and 66 are 2 standard deviations below and above 58,}\\
        \ansline{so about 95\% of the eggs (empirical rule).}

        \boxguard[5]
        \vspace{6pt}
        \textbf{(b)} About what percent of the eggs weigh \emph{more than} 62 grams?\\[4pt]
        \ansline{62 is 1 SD above. 68\% are within 1 SD, so 32\% are outside, split evenly:}\\
        \ansline{about 16\% of the eggs weigh more than 62 grams.}

        \boxguard[8]
        \vspace{6pt}
        \textbf{(c)} One egg weighs 51 grams. Find its $z$-score and interpret it in context.
        \begin{work}
          z &= \frac{51 - 58}{4} = -1.75
        \end{work}
        \ansline{This egg weighs 1.75 SDs below the mean egg weight of 58 grams.}

        \boxguard[10]
        \vspace{6pt}
        \textbf{(d)} A second farm's eggs have mean 64 grams and standard deviation 6 grams. Which
        egg is heavier \emph{relative to its own farm}: a 65-gram egg from the first farm or a
        73-gram egg from the second? Justify.
        \begin{work}
          z_{\text{first}} &= \frac{65 - 58}{4} = 1.75, \qquad
          z_{\text{second}} = \frac{73 - 64}{6} = 1.5
        \end{work}
        \ansline{The 65-gram egg: it is 1.75 SDs above its farm's mean, the 73-gram egg only}\\
        \ansline{1.5 --- relatively heavier even though it weighs fewer grams.}

\end{enumerate}

\end{document}
